



Ở mức thiết kế quan niệm (\textit{Conceptual Design Level}), mục tiêu cốt lõi là chuyển hóa các mô tả nghiệp vụ thô, yêu cầu chức năng và quy tắc hệ thống thành mô hình dữ liệu trừu tượng, hoàn toàn độc lập với Hệ quản trị cơ sở dữ liệu (DBMS). Đồ án áp dụng \textbf{Thuật toán 2.1} (\textit{Algorithm 2.1: Thiết kế Sơ đồ ER từ Yêu cầu Nghiệp vụ}) để đảm bảo tính toàn vẹn và phản ánh chính xác các nghiệp vụ cốt lõi của hệ thống \textbf{LivingDocs — Design and Implementation of an AI-Assisted Self-Updating Document System for Engineering Teams}.

Quá trình này tiếp nhận đầu vào (\textit{Input}) là bài toán giải quyết hiện tượng trôi lệch tài liệu (\textit{Documentation Drift}), bộ yêu cầu chức năng từ 5 vai trò người dùng (\textit{Developer, Staff, Manager, Technical Lead, Administrator}) và kết xuất đầu ra (\textit{Output}) là Sơ đồ Quan niệm (\textit{Conceptual ERD/UML}) đã qua kiểm chứng.





\section{Phân loại Tác nhân \& Phạm vi Nghiệp vụ }
\label{sec:actor-classification}
\begingroup
\setlength{\tabcolsep}{4pt} % Khoảng đệm giữa các cột
\small % Dùng cỡ chữ \small giúp nội dung lấp đầy bảng đẹp mắt hơn

\begin{longtable}{|
>{\raggedright\arraybackslash}p{0.13\textwidth}|
>{\raggedright\arraybackslash}p{0.17\textwidth}|
>{\raggedright\arraybackslash}p{0.35\textwidth}|
>{\raggedright\arraybackslash}p{0.13\textwidth}|
>{\raggedright\arraybackslash}p{0.13\textwidth}|}
\hline

\textbf{Tác nhân} &
\textbf{Mục tiêu} &
\textbf{Nhiệm vụ \& Quyền hạn chính} &
\textbf{Phạm vi dữ liệu} &
\textbf{Đặc điểm cốt lõi} \\
\hline
\endhead % Lặp lại dòng tiêu đề khi bảng ngắt sang trang sau

\textbf{Developer} &
Khởi tạo và sử dụng tài liệu phần mềm. &
Kết nối GitHub và chọn Repository vào Workspace; yêu cầu AI sinh tài liệu theo Template; xem Grounding Evidence; chỉnh sửa bản nháp và Submit; bật/tắt Auto-update; xem Version History và Drift Alert. &
Workspace và Repository được cấp quyền; không có quyền Approve hoặc Publish. &
Là tác nhân \textbf{duy nhất khởi tạo} vòng đời tài liệu, bắt đầu luồng \textbf{AI $\rightarrow$ Staff $\rightarrow$  Manangerr}. \\
\hline

\textbf{Staff} &
Thực hiện review kỹ thuật cấp 1. &
Xem Review Queue; kiểm tra Grounding Evidence, Source File, Code Entity, Commit, Code Diff và Change Log; kiểm tra Template; Comment, sửa lỗi nhỏ; Approve hoặc Reject và trả về Developer. &
Toàn bộ bối cảnh kỹ thuật của tài liệu đang review; không có quyền Publish hoặc Commit. &
Là \textbf{cổng lọc kỹ thuật}, phát hiện sai sót và AI hallucination trước khi chuyển đến Manager. \\
\hline

\textbf{Manager} &
Phê duyệt cuối và quyết định xuất bản tài liệu. &
Xem kết quả Staff review, Comment và Grounding Evidence; thực hiện Approve, Reject hoặc Rollback; quản lý danh mục tài liệu cần Manager approval; xem lịch sử Approval. &
Tài liệu đã qua Staff review và thuộc phạm vi quản lý. &
Là tác nhân \textbf{duy nhất} có quyền Publish hoặc Commit tài liệu vào Repository thật. \\
\hline

\textbf{Technical Lead} &
Quản trị chất lượng và chính sách tài liệu. &
Quản lý Documentation Collection và Ownership; thiết lập chính sách và ngưỡng phát hiện tài liệu lỗi thời ; quản lý Template, Section, Heading, Placeholder và ánh xạ Placeholder $\leftrightarrow$ Code Entity. &
Toàn bộ Workspace và Repository thuộc phạm vi quản lý. &
Tập trung vào \textbf{Policy, Threshold và Template}, không review từng tài liệu như Manager. \\
\hline

\textbf{Admin} &
Vận hành hạ tầng và quản trị hệ thống. &
Quản lý Tài khoản \& Phân quyền:  User, Role, Permission, Workspace và Repository và cấu hình tích hợp \& hạ tầng như GitHub App, OAuth, GitHub Actions, tích hợp GitLab, Jira, Slack; cấu hình ngôn ngữ và AI Model. &
Toàn bộ hệ thống ở cấp độ hạ tầng và kỹ thuật; không review hoặc approve tài liệu. &
Là tác nhân \textbf{duy nhất} quản lý cấu hình AI Model như LLM và Embedding, nằm ngoài luồng phê duyệt tài liệu. \\
\hline

\end{longtable}
\endgroup
\section{Chi tiết các Thực thể Dữ liệu Cốt lõi (Core Data Entities)}


\section{Chi tiết các Thực thể Dữ liệu}

% LaTeX does not define a fourth-level sectioning command by default.
\providecommand{\subsubsubsection}[1]{\paragraph{#1}}

\subsection{Nhóm Quản trị \& Kết nối}

\subsubsection{User (Thực thể mạnh)}

Người dùng cuối trong hệ thống (Developer, Staff, Manager, Tech Lead, Admin).

\subsubsubsection{Thuộc tính Phức hợp (Composite Attribute)}

\begin{itemize}
    \item \textbf{FullName:} Thuộc tính phức hợp gồm:
    \begin{itemize}
        \item FirstName
        \item LastName
    \end{itemize}
\end{itemize}

\subsubsubsection{Thuộc tính Đa trị (Multi-valued Attribute)}

\begin{itemize}
    \item \textbf{NotificationPreferences (Cài đặt thông báo):}
    Lưu nhiều giá trị cùng lúc cho một người dùng
    (ví dụ: chọn đồng thời cả \texttt{EMAIL} và \texttt{SYSTEM}).
\end{itemize}

\subsubsubsection{Thuộc tính Đơn (Simple / Atomic Attribute)}

\begin{itemize}
    \item \textbf{user\_id (PK):UUID} Định danh duy nhất cho từng người dùng.
    \item \textbf{email:} Địa chỉ email bắt buộc dùng để đăng nhập và nhận thông báo hệ thống.
    \item \textbf{password\_hash:} Mật khẩu đã mã hóa. Lưu giá trị khi đăng nhập bằng Email/Password và giữ \textbf{NULL} khi dùng thuần GitHub OAuth.
    \item \textbf{status:} Trạng thái tài khoản (ACTIVE, INACTIVE, BLOCKED).
    \item \textbf{created\_at:} Thời điểm khởi tạo tài khoản.
    \item \textbf{updated\_at:} Thời điểm cập nhật thông tin gần nhất.
    \item \textbf{deleted\_at:} Soft delete nhằm giữ nguyên lịch sử giao dịch (review, comment, approval) khi người dùng rời hệ thống.
\end{itemize}


\subsection{Workspace (Thực thể mạnh)}

Không gian làm việc chứa repository và tài liệu của một nhóm/dự án.

\subsubsection{Thuộc tính Đơn (Simple / Atomic Attribute)}

\begin{itemize}
    \item \textbf{workspace\_id (PK):UUID} Định danh duy nhất cho từng không gian làm việc.
    \item \textbf{name:} Tên không gian làm việc/dự án (VD: ``Payment Service Platform'').
    \item \textbf{description:} Mô tả chi tiết về phạm vi và mục tiêu của workspace.
    \item \textbf{status:} Trạng thái hoạt động của workspace (ACTIVE, ARCHIVED).
    \item \textbf{created\_at:} Thời điểm khởi tạo workspace.
    \item \textbf{updated\_at:} Thời điểm cập nhật thông tin workspace gần nhất.
    \item \textbf{deleted\_at:} Cờ xóa mềm (Soft delete) giữ nguyên dữ liệu báo cáo và audit trail khi ẩn workspace.
\end{itemize}


\subsection{Role (Thực thể mạnh)}

Vai trò của người dùng, được scoped theo Workspace.

\subsubsection{Thuộc tính Đơn (Simple / Atomic Attribute)}

\begin{itemize}
    \item \textbf{role\_id (PK):UUID} Định danh duy nhất cho vai trò.
    \item \textbf{role\_name:} Tên vai trò (VD: Developer, Staff, Manager, Tech Lead, Admin).
    \item \textbf{description:} Mô tả tóm tắt mục đích và quyền hạn chung của vai trò.
    \item \textbf{created\_at:} Thời điểm khởi tạo vai trò trong hệ thống.
    \item \textbf{updated\_at:} Thời điểm cập nhật vai trò gần nhất.
\end{itemize}


\subsubsection{Permission (Thực thể mạnh)}

Quyền hạn chi tiết trên hệ thống (VD: DOC\_APPROVE, TEMPLATE\_EDIT).

\subsubsubsection{Thuộc tính Đơn (Atomic Attribute)}

\begin{itemize}
    \item \textbf{permission\_id (PK):UUID} Định danh duy nhất cho quyền.
    \item \textbf{permission\_code:} Mã quyền thao tác hệ thống (VD: DOC\_APPROVE, TEMPLATE\_EDIT).
    \item \textbf{description:} Giải thích chi tiết giới hạn của thao tác này.
    \item \textbf{created\_at:} Mốc thời gian quyền được khởi tạo vào danh mục hệ thống.
\end{itemize}


\subsubsection{RolePermission (Thực thể liên kết)}

Bản ghi phân quyền chi tiết cho từng vai trò trong hệ thống.

\subsubsubsection{Thuộc tính Đơn (Atomic Attribute)}

\begin{itemize}
    \item \textbf{role\_permission\_id (PK):UUID} Định danh duy nhất cho bản ghi phân quyền (UUID).
    \item \textbf{role\_id (FK):UUID} Khóa ngoại trỏ đến vai trò được gán quyền (roles).
    \item \textbf{permission\_id (FK):UUID} Khóa ngoại trỏ đến quyền hạn cụ thể (permissions).
    \item \textbf{created\_at:} Mốc thời gian quyền này được gán cho vai trò.
\end{itemize}


\subsubsection{WorkspaceMember (Thực thể mối kết hợp)}

Ghi nhận sự tham gia của User trong Workspace và vai trò (Role) tương ứng.

\subsubsubsection{Thuộc tính Đơn (Atomic Attribute)}

\begin{itemize}
    \item \textbf{member\_id (PK):UUID} Định danh duy nhất cho bản ghi tham gia.
    \item \textbf{status:} Trạng thái thành viên (ACTIVE, INACTIVE, REMOVED).
    \item \textbf{joined\_at:} Thời điểm người dùng chính thức tham gia.
    \item \textbf{created\_at:} Thời điểm gửi lời mời / tạo bản ghi.
    \item \textbf{updated\_at:} Thời điểm cập nhật vai trò hoặc trạng thái gần nhất.
\end{itemize}


\subsubsection{GitHubConnection (Thực thể phụ thuộc)}

Lưu thông tin xác thực OAuth của User với GitHub.

\paragraph{Thuộc tính Khóa (Key Attribute)}

\begin{itemize}
    \item \textbf{connection\_id:} Định danh duy nhất cho bản ghi kết nối OAuth.
\end{itemize}

\paragraph{Thuộc tính Đơn / Nguyên tố (Simple / Atomic Attributes)}

\begin{itemize}
    \item \textbf{github\_user\_id:} ID tài khoản người dùng trên hệ thống GitHub để đối soát định danh.
    \item \textbf{github\_username:} Tên tài khoản GitHub (VD: octocat).
    \item \textbf{access\_token:} Chuỗi token ủy quyền dùng để gọi GitHub API.
    \item \textbf{created\_at:} Thời điểm người dùng thực hiện liên kết tài khoản.
    \item \textbf{updated\_at:} Thời điểm làm mới token hoặc cập nhật quyền hạn gần nhất.
\end{itemize}

\subsubsubsection{Thuộc tính Đa trị (Multi-valued Attribute)}

\begin{itemize}
    \item \textbf{scopes:} Lưu danh sách các phạm vi quyền OAuth được cấp phép cùng lúc (VD: repo, read:user, workflow).
\end{itemize}


\subsection{Nhóm Mã nguồn \& Cấu trúc Code}

\subsubsection{Repository (Thực thể mạnh)}

Kho chứa mã nguồn được đồng bộ vào hệ thống.

\subsubsubsection{Thuộc tính Đơn (Atomic Attribute)}

\begin{itemize}
    \item \textbf{repository\_id (PK):UUID} Định danh duy nhất cho kho mã nguồn.
    \item \textbf{name:} Tên kho mã nguồn (VD: payment-gateway).
    \item \textbf{url:} Đường dẫn gốc đến kho chứa trên VCS.
    \item \textbf{default\_branch:} Nhánh mặc định dùng để phân tích code (VD: main, master).
    \item \textbf{sync\_status:} Trạng thái đồng bộ dữ liệu (SYNCING, SUCCESS, FAILED).
    \item \textbf{last\_synced\_at:} Thời điểm đồng bộ mã nguồn thành công gần nhất.
    \item \textbf{created\_at:} Thời điểm tích hợp kho mã nguồn vào hệ thống.
    \item \textbf{updated\_at:} Thời điểm cập nhật trạng thái hoặc cấu hình kho gần nhất.
\end{itemize}


\subsubsection{Branch (Thực thể yếu -- phụ thuộc Repository)}

Nhánh phát triển trong repository.

\subsubsubsection{Thuộc tính Đơn (Simple / Atomic Attribute)}

\begin{itemize}
    \item \textbf{branch\_id (PK):UUID} Định danh duy nhất cho nhánh.
    \item \textbf{branch\_name:} Tên nhánh (VD: main, staging, feature/auth).
    \item \textbf{is\_default:} Cờ xác định nhánh mặc định của repo (true/false).
    \item \textbf{created\_at:} Thời điểm nhánh được nhận diện và lưu vào hệ thống.
    \item \textbf{updated\_at:} Thời điểm cập nhật commit hash hoặc trạng thái nhánh gần nhất.
\end{itemize}


\subsubsection{Commit (Thực thể yếu -- phụ thuộc Branch/Repository)}

Bản ghi một lần commit mã nguồn.

\subsubsubsection{Thuộc tính Phức hợp (Composite Attribute)}

\begin{itemize}
    \item \textbf{AuthorInfo:} Thông tin tác giả thực hiện commit, gồm:
    \begin{itemize}
        \item AuthorName
        \item AuthorEmail
    \end{itemize}
\end{itemize}

\subsubsubsection{Thuộc tính Đơn (Simple / Atomic Attribute)}

\begin{itemize}
    \item \textbf{commit\_id (PK):UUID} Định danh nội bộ cho bản ghi commit.
    \item \textbf{commit\_hash:} Mã băm SHA-1/SHA-256 duy nhất từ Git (VD: 7a8f9b2...).
    \item \textbf{message:} Nội dung mô tả của commit.
    \item \textbf{committed\_at:} Thời điểm commit được thực hiện trên VCS (thời gian nghiệp vụ).
    \item \textbf{created\_at:} Thời điểm bản ghi commit được đồng bộ vào LivingDocs (thời gian hệ thống).
\end{itemize}


\subsubsection{PullRequest (Thực thể yếu -- phụ thuộc Repository)}

Yêu cầu hợp nhất nhánh từ GitHub.

\subsubsubsection{Thuộc tính Đa trị (Multi-valued Attribute)}

\begin{itemize}
    \item \textbf{labels:} Danh sách nhãn phân loại được gắn trên PR
    (VD: bug, enhancement, documentation, api-change).
\end{itemize}

\subsubsubsection{Thuộc tính Đơn (Simple / Atomic Attribute)}

\begin{itemize}
    \item \textbf{pr\_id (PK):UUID} Định danh nội bộ của bản ghi PR.
    \item \textbf{pr\_number:} Số thứ tự PR trên nền tảng gốc (VD: \#102, \#123).
    \item \textbf{title:} Tiêu đề của Pull Request.
    \item \textbf{description:} Nội dung văn bản mô tả chi tiết PR, là dữ liệu đầu vào cho AI.
    \item \textbf{source\_branch:} Nhánh nguồn chứa mã nguồn thay đổi.
    \item \textbf{target\_branch:} Nhánh đích dự định hợp nhất vào (VD: main, develop).
    \item \textbf{status:} Trạng thái vòng đời của PR (OPEN, MERGED, CLOSED).
    \item \textbf{created\_at:} Thời điểm mở PR trên VCS.
    \item \textbf{updated\_at:} Thời điểm PR cập nhật thông tin/trạng thái gần nhất.
    \item \textbf{merged\_at:} Thời điểm PR chính thức được chấp nhận hợp nhất (Nullable).
\end{itemize}


\subsubsection{CodeDiff (Thực thể yếu -- phụ thuộc Commit/PR)}

Chi tiết thay đổi mã nguồn trong commit hoặc PR.

\subsubsubsection{Thuộc tính Đơn (Simple / Atomic Attribute)}

\begin{itemize}
    \item \textbf{diff\_id (PK):UUID} Định danh nội bộ bản ghi chi tiết thay đổi.
    \item \textbf{file\_path:} Đường dẫn tương đối đến tệp bị thay đổi (VD: src/controllers/payment.js).
    \item \textbf{change\_type:} Loại thao tác trên file (ADD, MODIFY, DELETE, RENAME).
    \item \textbf{additions:} Số dòng code được thêm mới.
    \item \textbf{deletions:} Số dòng code bị xóa đi.
    \item \textbf{diff\_content:} Nội dung chi tiết các đoạn code thay đổi (patch/hunk text) dùng làm input cho AI.
    \item \textbf{created\_at:} Thời điểm lưu bản ghi diff vào hệ thống.
\end{itemize}


\subsubsection{File (Thực thể yếu -- phụ thuộc Repository)}

Cấu trúc tập tin mã nguồn trong kho chứa.

\subsubsubsection{Thuộc tính Đơn (Simple / Atomic Attribute)}

\begin{itemize}
    \item \textbf{file\_id (PK):UUID} Định danh nội bộ của tập tin trong hệ thống.
    \item \textbf{file\_path:} Đường dẫn đầy đủ của file tính từ gốc repository.
    \item \textbf{file\_type:} Loại hoặc đuôi mở rộng của tệp (VD: ts, java, py, md).
    \item \textbf{status:} Trạng thái của file trong cây thư mục (ACTIVE, DELETED, RENAMED).
    \item \textbf{created\_at:} Thời điểm file lần đầu được ghi nhận vào hệ thống.
    \item \textbf{updated\_at:} Thời điểm thông tin file hoặc nội dung file thay đổi gần nhất.
\end{itemize}


\subsubsection{CodeEntity (Thực thể yếu -- phụ thuộc File)}

Đơn vị mã nguồn cụ thể thu được từ AST (Class, Method, Function...).

\subsubsubsection{Thuộc tính Phức hợp (Composite Attribute)}

\begin{itemize}
    \item \textbf{position:} Vị trí của thực thể trong tập tin, gồm:
    \begin{itemize}
        \item \textbf{start\_line:} Dòng bắt đầu của thực thể code.
        \item \textbf{end\_line:} Dòng kết thúc của thực thể code.
    \end{itemize}
\end{itemize}

\subsubsubsection{Thuộc tính Đơn (Simple / Atomic Attribute)}

\begin{itemize}
    \item \textbf{entity\_id (PK):UUID} Định danh nội bộ của thực thể code.
    \item \textbf{entity\_name:} Tên của thành phần code (VD: processPayment, UserController).
    \item \textbf{entity\_type:} Loại thành phần code (CLASS, METHOD, FUNCTION, INTERFACE, ENDPOINT).
    \item \textbf{signature:} Cấu trúc chữ ký của thành phần code.
    \item \textbf{doc\_comment:} Đoạn comment/docstring nguyên bản đính kèm hàm/lớp (nếu có).
    \item \textbf{created\_at:} Thời điểm thực thể code được bóc tách và lưu vào hệ thống.
    \item \textbf{updated\_at:} Thời điểm cấu trúc của thực thể code thay đổi gần nhất.
\end{itemize}


\subsubsection{Dependency (Thực thể mối kết hợp)}

Mối quan hệ phụ thuộc giữa các CodeEntity.

\subsubsubsection{Thuộc tính Đơn (Simple / Atomic Attribute)}

\begin{itemize}
    \item \textbf{dependency\_id (PK):UUID} Định danh duy nhất cho bản ghi phụ thuộc.
    \item \textbf{dependency\_type:} Loại quan hệ phụ thuộc (CALLS, INHERITS, IMPLEMENTS, USES, IMPORTS).
    \item \textbf{created\_at:} Thời điểm hệ thống phân tích AST và ghi nhận mối quan hệ phụ thuộc.
\end{itemize}


\subsection{Nhóm Khuôn mẫu Tài liệu (Template Engine Domain)}

\subsubsection{Template (Thực thể mạnh)}

Khuôn mẫu chuẩn để AI/User sinh tài liệu.

\subsubsubsection{Thuộc tính Đơn (Simple / Atomic Attribute)}

\begin{itemize}
    \item \textbf{template\_id (PK):UUID} Định danh duy nhất cho khuôn mẫu tài liệu.
    \item \textbf{name:} Tên khuôn mẫu.
    \item \textbf{description:} Mô tả mục đích sử dụng và phạm vi áp dụng.
    \item \textbf{content\_structure:} Nội dung cấu trúc dàn ý mẫu/Prompt Markdown hướng dẫn AI và User.
    \item \textbf{creator\_type:} Đối tượng khởi tạo khuôn mẫu (SYSTEM, USER, AI).
    \item \textbf{status:} Trạng thái sử dụng (DRAFT, ACTIVE, ARCHIVED).
    \item \textbf{created\_at:} Thời điểm tạo khuôn mẫu.
    \item \textbf{updated\_at:} Thời điểm cập nhật dàn ý/nội dung khuôn mẫu gần nhất.
\end{itemize}


\subsubsection{TemplateSection (Thực thể yếu -- phụ thuộc Template)}

Mục con/heading trong cấu trúc dàn ý của Template.

\begin{itemize}
    \item \textbf{section\_id (PK):UUID} Định danh duy nhất cho mục dàn ý.
    \item \textbf{heading\_title:} Tiêu đề của mục.
    \item \textbf{order\_index:} Thứ tự sắp xếp của mục trong Template.
    \item \textbf{prompt\_instruction:} Câu lệnh hướng dẫn riêng cho AI khi sinh nội dung cho mục này.
    \item \textbf{created\_at:} Thời điểm khởi tạo mục dàn ý.
    \item \textbf{updated\_at:} Thời điểm cập nhật tiêu đề hoặc chỉ thị gần nhất.
\end{itemize}


\subsubsection{Placeholder (Thực thể yếu -- phụ thuộc TemplateSection)}

Biến/điểm chèn trong Section, làm nhiệm vụ nhận dữ liệu từ code.

\begin{itemize}
    \item \textbf{placeholder\_id (PK):UUID} Định danh duy nhất cho điểm chèn biến.
    \item \textbf{variable\_name:} Mã nhận diện biến.
    \item \textbf{expected\_data\_type:} Loại dữ liệu kỳ vọng từ code.
    \item \textbf{description:} Mô tả ý nghĩa và ngữ cảnh dữ liệu cần điền vào biến.
    \item \textbf{is\_required:} Cờ xác định biến này bắt buộc AI phải lấp đầy hay không.
    \item \textbf{created\_at:} Thời điểm tạo điểm chèn biến.
    \item \textbf{updated\_at:} Thời điểm cập nhật thông tin biến gần nhất.
\end{itemize}


\subsubsection{TemplateMapping (Thực thể mối kết hợp)}

Ánh xạ trực tiếp giữa Placeholder và CodeEntity để AI trích xuất thông tin.

\begin{itemize}
    \item \textbf{mapping\_id (PK):UUID} Định danh duy nhất cho bản ghi ánh xạ trích xuất.
    \item \textbf{extraction\_rule:} Quy tắc/Prompt chỉ thị chi tiết để AI biết cách lọc, biến đổi và trích xuất dữ liệu từ CodeEntity.
    \item \textbf{is\_active:} Trạng thái hiệu lực của quy tắc ánh xạ.
    \item \textbf{created\_at:} Thời điểm khởi tạo quy tắc ánh xạ.
    \item \textbf{updated\_at:} Thời điểm cập nhật quy tắc trích xuất gần nhất.
\end{itemize}


\subsection{Nhóm Tài liệu \& Vòng đời Dữ liệu (Documentation Domain)}

\subsubsection{Document (Thực thể mạnh)}

Tài liệu gốc (README, API Reference, Architecture Guide, ADR...).

\begin{itemize}
    \item \textbf{document\_id (PK):UUID} Định danh duy nhất cho tài liệu chính thức.
    \item \textbf{title:} Tên/Tiêu đề của tài liệu.
    \item \textbf{document\_type:} Phân loại loại tài liệu (README, API\_REFERENCE, ARCHITECTURE\_GUIDE, ADR, CUSTOM).
    \item \textbf{status:} Trạng thái vòng đời của tài liệu (PUBLISHED, ARCHIVED, DEPRECATED).
    \item \textbf{created\_at:} Thời điểm khởi tạo tài liệu.
    \item \textbf{updated\_at:} Thời điểm cập nhật tài liệu gần nhất.
    \item \textbf{deleted\_at:} Thời điểm xóa mềm tài liệu (Nullable).
\end{itemize}


\subsubsection{DocumentVersion (Thực thể yếu -- phụ thuộc Document)}

Bản ghi nội dung tài liệu theo từng mốc thời gian và trạng thái.

\begin{itemize}
    \item \textbf{version\_id (PK):UUID} Định danh duy nhất cho bản ghi phiên bản.
    \item \textbf{version\_number:} Số hiệu phiên bản tài liệu (VD: 1.0.0, 1.1.0, 2.0-draft).
    \item \textbf{content:} Toàn bộ nội dung văn bản tài liệu dạng Markdown/HTML tại thời điểm chụp phiên bản.
    \item \textbf{change\_summary:} Tóm tắt mô tả lý do hoặc nội dung những điểm thay đổi trong phiên bản này.
    \item \textbf{status:} Trạng thái vòng đời phiên bản (DRAFT, IN\_REVIEW, APPROVED, PUBLISHED).
    \item \textbf{ticket\_codes:} Danh sách các mã ticket/thẻ công việc liên quan từ các hệ thống bên ngoài như Jira Ticket ID hoặc GitHub Issue ID. Thuộc tính này được lưu dưới dạng chuỗi văn bản đơn giản, Nullable và phân cách bằng dấu phẩy.
    \item \textbf{created\_at:} Thời điểm khởi tạo bản ghi phiên bản.
    \item \textbf{updated\_at:} Thời điểm cập nhật trạng thái phiên bản gần nhất.
\end{itemize}


\subsubsection{ChangeLog (Thực thể yếu -- phụ thuộc DocumentVersion)}

Lịch sử vết thay đổi chi tiết giữa các phiên bản tài liệu.

\begin{itemize}
    \item \textbf{log\_id (PK):UUID} Định danh duy nhất cho bản ghi nhật ký thay đổi.
    \item \textbf{change\_type:} Loại thao tác biến đổi (ADDED, MODIFIED, DELETED, RESTORED).
    \item \textbf{change\_description:} Mô tả tóm tắt ý nghĩa sự thay đổi.
    \item \textbf{diff\_content:} Nội dung chênh lệch chi tiết (dạng Git patch / JSON delta).
    \item \textbf{created\_at:} Thời điểm ghi nhận vết thay đổi.
    \item \textbf{updated\_at:} Thời điểm cập nhật bản ghi nhật ký nếu có bổ sung metadata.
\end{itemize}


\subsection{Nhóm Kiểm duyệt \& Phê duyệt (Review \& Governance Domain)}

\subsubsection{Review (Thực thể mối kết hợp/Sự kiện)}

Tiến trình Staff thực hiện đánh giá kỹ thuật trên một DocumentVersion.

\begin{itemize}
    \item \textbf{review\_id (PK):UUID} Định danh duy nhất cho phiên đánh giá kỹ thuật.
    \item \textbf{status:} Trạng thái của tiến trình review (PENDING, IN\_PROGRESS, COMPLETED).
    \item \textbf{decision:} Kết quả thẩm định của reviewer (APPROVED, REJECTED, CHANGES\_REQUESTED).
    \item \textbf{summary\_notes:} Nhận xét hoặc ghi chú tổng quan của reviewer (Nullable).
    \item \textbf{created\_at:} Thời điểm bắt đầu phiên đánh giá.
    \item \textbf{updated\_at:} Thời điểm chốt kết quả hoặc cập nhật tiến trình review.
\end{itemize}


\subsubsection{ReviewComment (Thực thể yếu -- phụ thuộc Review)}

Phản hồi, góp ý chi tiết trong quá trình kiểm duyệt.

\begin{itemize}
    \item \textbf{comment\_id (PK):UUID} Định danh duy nhất cho bản ghi bình luận/góp ý.
    \item \textbf{content:} Nội dung chi tiết của lời nhận xét, phản hồi.
    \item \textbf{is\_resolved:} Trạng thái xử lý góp ý (false: Chưa sửa / true: Đã sửa xong).
    \item \textbf{created\_at:} Thời điểm gửi bình luận.
    \item \textbf{updated\_at:} Thời điểm cập nhật nội dung bình luận gần nhất.
\end{itemize}


\subsubsection{Approval (Thực thể mối kết hợp/Sự kiện)}

Quyết định phê duyệt xuất bản/commit tài liệu chính thức từ Manager.

\begin{itemize}
    \item \textbf{approval\_id (PK):UUID} Định danh duy nhất cho bản ghi quyết định phê duyệt.
    \item \textbf{decision:} Trạng thái/Kết quả quyết định (APPROVED, REJECTED).
    \item \textbf{rejection\_reason:} Lý do từ chối xuất bản (Nullable, dùng khi decision = REJECTED).
    \item \textbf{approved\_at:} Mốc thời gian ra quyết định phê duyệt hoặc từ chối.
\end{itemize}


\subsection{Nhóm AI Evidence, Giám sát \& Kiểm toán}

\subsubsection{GroundingEvidence (Thực thể yếu -- phụ thuộc DocumentVersion)}

Bảng minh chứng chứa liên kết nguồn gốc (File, CodeEntity, Commit) AI dùng để tạo nội dung, phục vụ chống hallucination.

\begin{itemize}
    \item \textbf{evidence\_id (PK):UUID} Định danh duy nhất cho bản ghi bằng chứng minh chứng.
    \item \textbf{snippet\_content:} Đoạn mã nguồn hoặc văn bản thô trích xuất từ code làm căn cứ trực tiếp cho AI sinh nội dung.
    \item \textbf{purpose:} Phân loại mục đích lưu vết (HALLUCINATION\_PREVENTION, CODE\_REFERENCE, AST\_VERIFICATION).
    \item \textbf{created\_at:} Thời điểm ghi nhận bản ghi bằng chứng.
\end{itemize}


\subsubsection{DriftAlert (Thực thể sự kiện)}

Cảnh báo phát hiện tài liệu bị lệch ngữ nghĩa/lỗi thời so với code mới.

\subsubsubsection{Thuộc tính Khóa (Key Attribute)}

\begin{itemize}
    \item \textbf{alert\_id:} Định danh duy nhất cho từng bản ghi cảnh báo.
\end{itemize}

\subsubsubsection{Thuộc tính Đơn / Nguyên tố (Simple / Atomic Attributes)}

\begin{itemize}
    \item \textbf{drift\_type:} Phân loại hình thức lệch pha (SEMANTIC, OUTDATED).
    \item \textbf{status:} Trạng thái xử lý cảnh báo (OPEN, RESOLVED, IGNORED).
    \item \textbf{detected\_at:} Thời điểm hệ thống AI quét và phát hiện ra độ lệch.
\end{itemize}

\subsection{Thực thể mối kết hợp / Sự kiện: Approval}

\begin{itemize}
    \item \textbf{Tên thực thể:} \texttt{Approval} (\textit{Phê duyệt})
    \item \textbf{Mô tả nghiệp vụ:} Quyết định phê duyệt xuất bản hoặc \textit{commit} tài liệu chính thức từ phía Manager.
    \item \textbf{Thuộc tính đơn \& Khóa (Atomic \& Key Attributes):}
    \begin{itemize}
        \item \texttt{approval\_id} (\textbf{PK}): Định danh duy nhất cho bản ghi quyết định phê duyệt (\texttt{UUID}).
        \item \texttt{version\_id} (\textbf{FK}): Khóa ngoại trỏ đến phiên bản tài liệu (\texttt{document\_versions}) được xét duyệt.
        \item \texttt{user\_id} (\textbf{FK}): Khóa ngoại trỏ đến Manager (\texttt{users}) thực hiện ra quyết định.
        \item \texttt{decision}: Trạng thái/Kết quả quyết định (\texttt{APPROVED}, \texttt{REJECTED}).
        \item \texttt{rejection\_reason}: Lý do từ chối xuất bản (\textit{Nullable}, dùng khi \texttt{decision} = \texttt{REJECTED}).
        \item \texttt{approved\_at}: Mốc thời gian ra quyết định phê duyệt hoặc từ chối (\texttt{DATETIME2} / \texttt{TIMESTAMP}).
    \end{itemize}
\end{itemize}
\section{Mô hình hóa Mối kết hợp \& Bản số Dữ liệu}

\begingroup
\setlength{\tabcolsep}{3.5pt} % Thu nhỏ khoảng đệm giữa các cột
\small % Cỡ chữ vừa vặn giúp bảng đẹp và cân đối

\begin{longtable}{|
>{\raggedright\arraybackslash}p{0.24\textwidth}|
>{\raggedright\arraybackslash}p{0.25\textwidth}|
>{\raggedright\arraybackslash}p{0.17\textwidth}|
>{\raggedright\arraybackslash}p{0.34\textwidth}|}
\hline

\textbf{Mối quan hệ \& Phân loại thực thể} &
\textbf{Bản số \& Mức tham gia} &
\textbf{Thuộc tính mối quan hệ} &
\textbf{Mô tả nghiệp vụ} \\
\hline
\endfirsthead

\hline
\textbf{Mối quan hệ \& Phân loại thực thể} &
\textbf{Bản số \& Mức tham gia} &
\textbf{Thuộc tính mối quan hệ} &
\textbf{Mô tả nghiệp vụ} \\
\hline
\endhead

\textbf{User — Workspace — Role} (\texttt{<participates>}) \newline
\textit{Mối kết hợp 3 ngôi:} User (Mạnh), Workspace (Mạnh), Role (Mạnh) &
\textbf{Tỷ lệ:} M:N:1 \newline
\textbf{Bản số:} User (0..*), Workspace (1..*), Role (0..*) \newline
\textbf{Tham gia:} User (Partial 0), Workspace (Total 1), Role (Partial 0) &
\texttt{joined\allowbreak\_at}, \texttt{status} &
Ghi nhận sự tham gia của Người dùng vào Workspace. Trong từng Workspace, mỗi user bắt buộc gán đúng một Role (Scoped RBAC). Một user có thể thuộc nhiều Workspace với vai trò độc lập. \\ \hline

\textbf{Role — Permission} (\texttt{<contains>}) \newline
Role (Mạnh) $\leftrightarrow$ Permission (Mạnh) &
\textbf{Tỷ lệ:} M:N \newline
\textbf{Bản số:} Role (0..*) $\leftrightarrow$ Permission (0..*) \newline
\textbf{Tham gia:} Role (Partial) $\leftrightarrow$ Permission (Partial) &
\texttt{granted\allowbreak\_at} &
Một vai trò chứa tập hợp nhiều quyền; một quyền hạn có thể thuộc về nhiều vai trò khác nhau. \\ \hline

\textbf{User — Github\allowbreak\_Connection} (\texttt{<connects>}) \newline
User (Mạnh) $\leftrightarrow$ Github\_Connection (Phụ thuộc User) &
\textbf{Tỷ lệ:} 1:N \newline
\textbf{Bản số:} User (0..*), Github\allowbreak\_Connection (1..1) \newline
\textbf{Tham gia:} User (Partial), Github\allowbreak\_Connection (Total) &
Không có &
Ghi nhận việc User ủy quyền tài khoản GitHub qua OAuth2, lưu \texttt{access\_token} và scopes để đồng bộ repo, commit, PR, webhook. \\ \hline

\textbf{Workspace — Repository} (\texttt{<owns>}) \newline
Workspace (Mạnh) $\leftrightarrow$ Repository (Mạnh) &
\textbf{Tỷ lệ:} 1:N \newline
\textbf{Bản số:} Workspace (1..1) $\leftrightarrow$ Repository (0..*) \newline
\textbf{Tham gia:} Workspace (Partial), Repository (Total) &
Không có &
Một Workspace quản lý nhiều Repository; mỗi Repository thuộc quyền sở hữu của một Workspace. \\ \hline

\textbf{Repository — Branch} (\texttt{<has\_branch>}) \newline
Repository (Mạnh) $\leftrightarrow$ Branch (Yếu - Phụ thuộc Repo) &
\textbf{Tỷ lệ:} 1:N \newline
\textbf{Bản số:} Repository (1..1) $\leftrightarrow$ Branch (1..*) \newline
\textbf{Tham gia:} Repository (Total), Branch (Total) &
Không có &
Một kho chứa mã nguồn có ít nhất một nhánh mặc định; mỗi nhánh bắt buộc nằm trong duy nhất một kho chứa. \\ \hline

\textbf{Branch — Commit} (\texttt{<contains\_commit>}) \newline
Branch (Yếu) $\leftrightarrow$ Commit (Yếu - Phụ thuộc Branch/Repo) &
\textbf{Tỷ lệ:} 1:N \newline
\textbf{Bản số:} Branch (0..*) $\leftrightarrow$ Commit (1..1) \newline
\textbf{Tham gia:} Branch (Partial), Commit (Total) &
Không có &
Một nhánh chứa nhiều lượt commit; mỗi commit thuộc về đúng một nhánh phát triển. \\ \hline

\textbf{User — Commit} (\texttt{<makes\_commit>}) \newline
User (Mạnh) $\leftrightarrow$ Commit (Yếu) &
\textbf{Tỷ lệ:} 1:N \newline
\textbf{Bản số:} User (0..*) $\leftrightarrow$ Commit (1..1) \newline
\textbf{Tham gia:} User (Partial), Commit (Total) &
Không có &
Một người dùng thực hiện nhiều commit; mỗi commit được tạo bởi đúng một người dùng. \\ \hline

\textbf{Repository — Pull\allowbreak Request} (\texttt{<receives\_pr>}) \newline
Repository (Mạnh) $\leftrightarrow$ PullRequest (Yếu - Phụ thuộc Repo) &
\textbf{Tỷ lệ:} 1:N \newline
\textbf{Bản số:} Repository (0..*) $\leftrightarrow$ PullRequest (1..1) \newline
\textbf{Tham gia:} Repository (Partial), PullRequest (Total) &
Không có &
Một Repository chứa nhiều Yêu cầu hợp nhất (PR); mỗi PR gửi về đúng một kho chứa. \\ \hline

\textbf{Commit — Code\allowbreak Diff} (\texttt{<generates\_diff>}) \newline
Commit (Yếu) $\leftrightarrow$ CodeDiff (Yếu - Phụ thuộc Commit/PR) &
\textbf{Tỷ lệ:} 1:N \newline
\textbf{Bản số:} Commit (1..*) $\leftrightarrow$ CodeDiff (1..1) \newline
\textbf{Tham gia:} Commit (Total), CodeDiff (Total) &
Không có &
Một Commit sinh ra một hoặc nhiều bản ghi khác biệt code (Diff); mỗi bản ghi diff thuộc về duy nhất một commit. \\ \hline

\textbf{Pull\allowbreak Request — Commit} (\texttt{<includes\_commit>}) \newline
PullRequest (Yếu) $\leftrightarrow$ Commit (Yếu) &
\textbf{Tỷ lệ:} 1:N \newline
\textbf{Bản số:} PullRequest (1..*) $\leftrightarrow$ Commit (1..1) \newline
\textbf{Tham gia:} PullRequest (Total), Commit (Partial) &
Không có &
Một Yêu cầu hợp nhất (PR) tổng hợp danh sách gồm một hoặc nhiều commit; mỗi commit thuộc về đúng một PR. \\ \hline

\textbf{Commit — File} (\texttt{<modifies>}) \newline
Commit (Yếu) $\leftrightarrow$ File (Yếu - Phụ thuộc Repo) &
\textbf{Tỷ lệ:} M:N \newline
\textbf{Bản số:} Commit (1..*) $\leftrightarrow$ File (0..*) \newline
\textbf{Tham gia:} Commit (Total), File (Partial) &
\texttt{change\allowbreak\_type}, \newline
\texttt{lines\allowbreak\_added}, \newline
\texttt{lines\allowbreak\_deleted} &
Một Commit tác động chỉnh sửa ít nhất một File; một File có thể chịu sự biến đổi qua nhiều lượt Commit. \\ \hline

\textbf{Repository — File} (\texttt{<contains\_file>}) \newline
Repository (Mạnh) $\leftrightarrow$ File (Yếu - Phụ thuộc Repo) &
\textbf{Tỷ lệ:} 1:N \newline
\textbf{Bản số:} Repository (0..*) $\leftrightarrow$ File (1..1) \newline
\textbf{Tham gia:} Repository (Partial), File (Total) &
Không có &
Một Repository quản lý cấu trúc cây thư mục chứa nhiều File mã nguồn; mỗi File thuộc một kho chứa duy nhất. \\ \hline

\textbf{File — Code\allowbreak Entity} (\texttt{<parses\_into>}) \newline
File (Yếu) $\leftrightarrow$ CodeEntity (Yếu - Phụ thuộc File) &
\textbf{Tỷ lệ:} 1:N \newline
\textbf{Bản số:} File (0..*) $\leftrightarrow$ Code\allowbreak Entity (1..1) \newline
\textbf{Tham gia:} File (Partial), Code\allowbreak Entity (Total) &
Không có &
Một File code qua bóc tách AST thu được nhiều thực thể code (Class, Function); mỗi thực thể code nằm trọn trong đúng một File. \\ \hline

\textbf{Code\allowbreak Entity — Code\allowbreak Entity} (\texttt{<depends\_on>}) \newline
\textit{Đệ quy:} CodeEntity (Yếu) $\leftrightarrow$ CodeEntity (Yếu) &
\textbf{Tỷ lệ:} M:N \newline
\textbf{Bản số:} Phụ thuộc (0..*), Được phụ thuộc (0..*) \newline
\textbf{Tham gia:} Cả 2 phía Partial &
\texttt{dependency\allowbreak\_type} &
Xác định mối quan hệ phụ thuộc (gọi hàm, kế thừa, import) giữa các đơn vị code trong hệ thống. \\ \hline

\textbf{Workspace — Template} (\texttt{<defines>}) \newline
Workspace (Mạnh) $\leftrightarrow$ Template (Mạnh) &
\textbf{Tỷ lệ:} 1:N \newline
\textbf{Bản số:} Workspace (0..*) $\leftrightarrow$ Template (1..1) \newline
\textbf{Tham gia:} Workspace (Partial), Template (Total) &
Không có &
Một Workspace định nghĩa nhiều khuôn mẫu tài liệu; mỗi mẫu tài liệu thuộc sở hữu của một Workspace. \\ \hline

\textbf{Template — Template\allowbreak Section} (\texttt{<has\_section>}) \newline
Template (Mạnh) $\leftrightarrow$ TemplateSection (Yếu - Phụ thuộc Template) &
\textbf{Tỷ lệ:} 1:N \newline
\textbf{Bản số:} Template (1..*) $\leftrightarrow$ Template\allowbreak Section (1..1) \newline
\textbf{Tham gia:} Template (Total), Template\allowbreak Section (Total) &
Không có &
Một Template cấu thành từ ít nhất một Section dàn ý; mỗi Section nằm trên đúng một Template. \\ \hline

\textbf{Template\allowbreak Section — Placeholder} (\texttt{<contains\_placeholder>}) \newline
TemplateSection (Yếu) $\leftrightarrow$ Placeholder (Yếu - Phụ thuộc Section) &
\textbf{Tỷ lệ:} 1:N \newline
\textbf{Bản số:} Template\allowbreak Section (0..*) $\leftrightarrow$ Placeholder (1..1) \newline
\textbf{Tham gia:} Template\allowbreak Section (Partial), Placeholder (Total) &
Không có &
Mỗi mục trong mẫu tài liệu chứa nhiều biến chèn tự động. Khi xóa Section, các Placeholder nằm trong đó bị xóa theo. \\ \hline

\textbf{Placeholder — Code\allowbreak Entity} (\texttt{<maps\_to>}) \newline
Placeholder (Yếu) $\leftrightarrow$ CodeEntity (Yếu) &
\textbf{Tỷ lệ:} M:N \newline
\textbf{Bản số:} Placeholder (0..*) $\leftrightarrow$ Code\allowbreak Entity (0..*) \newline
\textbf{Tham gia:} Placeholder (Partial), Code\allowbreak Entity (Partial) &
\texttt{extraction\allowbreak\_rule}, \newline
\texttt{is\allowbreak\_active} &
Ánh xạ chỉ thị cho AI biết lấy thông tin từ thành phần code nào để điền vào biến trong mẫu. \\ \hline

\textbf{Workspace — Document} (\texttt{<manages>}) \newline
Workspace (Mạnh) $\leftrightarrow$ Document (Mạnh) &
\textbf{Tỷ lệ:} 1:N \newline
\textbf{Bản số:} Workspace (0..*) $\leftrightarrow$ Document (1..1) \newline
\textbf{Tham gia:} Workspace (Partial), Document (Total) &
Không có &
Workspace quản lý danh mục tài liệu gốc của dự án; một tài liệu thuộc về đúng một Workspace. \\ \hline

\textbf{User — Document} (\texttt{<authors>}) \newline
User (Mạnh) $\leftrightarrow$ Document (Mạnh) &
\textbf{Tỷ lệ:} 1:N \newline
\textbf{Bản số:} User (0..*) $\leftrightarrow$ Document (1..1) \newline
\textbf{Tham gia:} User (Partial), Document (Total) &
Không có &
Người dùng đứng tên tác giả khởi tạo tài liệu; mỗi tài liệu có một tác giả duy nhất. \\ \hline

\textbf{Document — Document\allowbreak Version} (\texttt{<has\_version>}) \newline
Document (Mạnh) $\leftrightarrow$ DocumentVersion (Yếu - Phụ thuộc Document) &
\textbf{Tỷ lệ:} 1:N \newline
\textbf{Bản số:} Document (1..*) $\leftrightarrow$ Document\allowbreak Version (1..1) \newline
\textbf{Tham gia:} Document (Total), Document\allowbreak Version (Total) &
Không có &
Một tài liệu chính thức chứa ít nhất một phiên bản nội dung; mỗi phiên bản đại diện cho một mốc chỉnh sửa. \\ \hline

\textbf{Document\allowbreak Version — Change\allowbreak Log} (\texttt{<tracks\_change>}) \newline
DocumentVersion (Yếu) $\leftrightarrow$ ChangeLog (Yếu - Phụ thuộc DocVersion) &
\textbf{Tỷ lệ:} 1:N \newline
\textbf{Bản số:} Document\allowbreak Version (0..*) $\leftrightarrow$ ChangeLog (1..1) \newline
\textbf{Tham gia:} Document\allowbreak Version (Partial), ChangeLog (Total) &
Không có &
Mỗi phiên bản tài liệu có thể ghi nhận nhiều lượt nhật ký biến động nội dung chi tiết. \\ \hline

\textbf{Document\allowbreak Version — Review} (\texttt{<has\_review>}) \newline
DocumentVersion (Yếu) $\leftrightarrow$ Review (Yếu - Phụ thuộc DocVersion) &
\textbf{Tỷ lệ:} 1:N \newline
\textbf{Bản số:} Document\allowbreak Version (0..*) $\leftrightarrow$ Review (1..1) \newline
\textbf{Tham gia:} Document\allowbreak Version (Partial), Review (Total) &
Không có &
Một phiên bản tài liệu có thể nhận 0 hoặc nhiều lượt đánh giá kỹ thuật; mỗi lượt đánh giá thuộc về duy nhất một phiên bản. \\ \hline

\textbf{User — Review} (\texttt{<performs\_review>}) \newline
User (Mạnh) $\leftrightarrow$ Review (Yếu) &
\textbf{Tỷ lệ:} 1:N \newline
\textbf{Bản số:} User (0..*) $\leftrightarrow$ Review (1..1) \newline
\textbf{Tham gia:} User (Partial), Review (Total) &
Không có &
Một người dùng có thể thực hiện 0 hoặc nhiều lượt đánh giá; mỗi lượt review do đúng một người dùng chịu trách nhiệm. \\ \hline

\textbf{Review — Review\allowbreak Comment} (\texttt{<contains\_comment>}) \newline
Review (Sự kiện/Yếu) $\leftrightarrow$ ReviewComment (Yếu - Phụ thuộc Review) &
\textbf{Tỷ lệ:} 1:N \newline
\textbf{Bản số:} Review (0..*) $\leftrightarrow$ Review\allowbreak Comment (1..1) \newline
\textbf{Tham gia:} Review (Partial), Review\allowbreak Comment (Total) &
Không có &
Một phiên kiểm duyệt kỹ thuật chứa nhiều bình luận góp ý chi tiết; mỗi bình luận gắn chặt với một lượt review. \\ \hline

\textbf{Document\allowbreak Version — Approval} (\texttt{<requires\_approval>}) \newline
DocumentVersion (Yếu) $\leftrightarrow$ Approval (Yếu - Phụ thuộc DocVersion) &
\textbf{Tỷ lệ:} 1:N \newline
\textbf{Bản số:} Document\allowbreak Version (0..*) $\leftrightarrow$ Approval (1..1) \newline
\textbf{Tham gia:} Document\allowbreak Version (Partial), Approval (Total) &
Không có &
Một phiên bản tài liệu có thể có 0 hoặc nhiều lượt phê duyệt; mỗi lượt phê duyệt bắt buộc thuộc về đúng một phiên bản. \\ \hline

\textbf{User — Approval} (\texttt{<grants\_approval>}) \newline
User (Mạnh) $\leftrightarrow$ Approval (Yếu) &
\textbf{Tỷ lệ:} 1:N \newline
\textbf{Bản số:} User (0..*) $\leftrightarrow$ Approval (1..1) \newline
\textbf{Tham gia:} User (Partial), Approval (Total) &
Không có &
Một Manager có thể thực hiện 0 hoặc nhiều đợt phê duyệt; mỗi bản ghi phê duyệt do đúng một Manager ra quyết định. \\ \hline

\textbf{Document\allowbreak Version — Ground\allowbreak ing\allowbreak Evidence} (\texttt{<has\_evidence>}) \newline
DocumentVersion (Yếu) $\leftrightarrow$ GroundingEvidence (Yếu - Phụ thuộc DocVersion) &
\textbf{Tỷ lệ:} 1:N \newline
\textbf{Bản số:} Document\allowbreak Version (0..*) $\leftrightarrow$ Grounding\allowbreak Evidence (1..1) \newline
\textbf{Tham gia:} Document\allowbreak Version (Partial), Grounding\allowbreak Evidence (Total) &
Không có &
Phiên bản tài liệu lưu trữ bảng minh chứng đối soát nguồn gốc code để chống ảo giác (hallucination) của AI. \\ \hline

\textbf{Ground\allowbreak ing\allowbreak Evidence — Code\allowbreak Entity} (\texttt{<references\_code>}) \newline
GroundingEvidence (Yếu) $\leftrightarrow$ CodeEntity (Yếu) &
\textbf{Tỷ lệ:} M:N \newline
\textbf{Bản số:} Grounding\allowbreak Evidence (0..*) $\leftrightarrow$ Code\allowbreak Entity (0..*) \newline
\textbf{Tham gia:} Grounding\allowbreak Evidence (Partial), Code\allowbreak Entity (Partial) &
\texttt{similarity\allowbreak\_score}, \newline
\texttt{match\allowbreak\_type} &
Minh chứng AI liên kết trực tiếp với các đơn vị code cụ thể trong hệ thống để phục vụ đối soát sai lệch. \\ \hline

\textbf{Ground\allowbreak ing\allowbreak Evidence — Drift\allowbreak Alert} (\texttt{<detects>}) \newline
GroundingEvidence (Yếu) $\leftrightarrow$ DriftAlert (Thực thể sự kiện) &
\textbf{Tỷ lệ:} 1:N \newline
\textbf{Bản số:} Grounding\allowbreak Evidence (0..*) $\leftrightarrow$ DriftAlert (1..1) \newline
\textbf{Tham gia:} Grounding\allowbreak Evidence (Partial), DriftAlert (Total) &
Không có &
Minh chứng đối soát khi phát hiện sự bất đồng giữa tài liệu và mã nguồn mới sẽ kích hoạt cảnh báo sai lệch (Drift Alert). \\ \hline

\end{longtable}
\endgroup


